% 16_body.tex — Day4 산출물 1/5 (⑯) 컷오버 런북 · 서비스 수용 기준(SAC) (빈 템플릿 / 완성 예시 공용 본문)
% 드라이버: 16_컷오버런북_SAC_빈템플릿.tex (\wbblanktrue) / 16_컷오버런북_SAC_완성예시.tex (\wbblankfalse)
% 내용 출처: PM트랙/Day4_워크북.md §1.3(항목 가이드) §1.4(빈 템플릿) §1.5(온담 P1 완성 예시본)
\documentclass[10pt,a4paper]{article}
\usepackage[a4paper,margin=16mm,top=14mm,bottom=20mm]{geometry}
\def\DocNum{1}
\def\DocTitle{컷오버 런북 · 서비스 수용 기준}
\def\DocSub{⑯ Cutover Runbook \& Service Acceptance Criteria}
\def\DocVariant{\B{빈 템플릿}{온담 P1 완성 예시}}
\usepackage{wbstyle}
\begin{document}

\docheader
 {\Bg{[정식 명칭과 코드 — 몇 차 컷오버인가]}{차세대 주문·재고 통합 플랫폼 구축 (ONE ONDAM P1) — 1차 컷오버(매장·주문)}}
 {\Bg{[v / 확정일 / 실행 기록일]}{v1.0 2027-02-05 확정 · v1.1 02-25 / 실행 기록 02-28 / 인수 서명 03-01}}
 {\Bg{[통제관 실명]}{P1 PM (통제관)\rd{7mm}}}
 {\Bg{[go/no-go 판정자 실명]}{정미란 CFO (게이트 3·오픈)\rd{7mm}}}

% ------------------------------------------------------------------ 1
\secbar{1. 문서 정보 — 창 · 시각 기준 · 역할 · 판정 권한 · 리허설}
\guide{컷오버 창(시작·종료·총 시간)과 T0, 버전과 확정일(리허설 몇 회차를 반영했는가), 컨트롤룸, 역할별 \textbf{실명}(통제관·데이터·인터페이스·매장 배포·검증·롤백·운영 인수·기록), 게이트별 판정자와 위임·\textbf{판정자 부재 시 대체자}, 연락망, 리허설 이력(회차·일자·소요·결과). 시각을 “1일차 오전”으로 적지 않는다. 판정자를 “PM”으로 적어 놓고 실제로는 회의로 정하지 않는다.}
{\footnotesize
\begin{tabularx}{\linewidth}{|L{28mm}|Y|}\hline
\hd{항목} & \hd{내용} \\ \hline
\B{\rd{9mm}}{}창 / T0 & \Bg{[YYYY-MM-DD(요일) hh:mm ~ YYYY-MM-DD hh:mm, n시간 / T0 = 창 시작. 헌장의 컷오버 창 조항·허용범위]}{2027-02-26(금) 20:00 ~ 02-28(일) 06:00, \textbf{34시간}(P-35) / T0 = 02-26 20:00. 헌장 §7 컷오버 창, 허용범위 0} \\ \hline
\B{\rd{9mm}}{}버전 & \Bg{[v1.0 확정일(반영한 리허설 회차) → v1.1 (변경 사유)]}{v1.0 02-05(리허설 3회차 01-22 실측 반영, G3 자료 첨부) → v1.1 02-25(항목 9 D$-$2 사건: 게이트 1 조건에 “임시 라이선스 유효” 추가)} \\ \hline
\B{\rd{8mm}}{}컨트롤룸 & \Bg{[장소 + 워룸 + 원격 상황실]}{본사 7층 대회의실(컨트롤룸) + 6층(수행사 워룸 20명) + 이천센터 상황실(한동석 대리, 폴백 배치 확인)} \\ \hline
\B{\rd{16mm}}{}역할 (실명) & \Bg{[통제관 / 데이터 / 인터페이스 / 매장 배포 / 검증 / 롤백 / 운영 인수 / 규제 / 기록 / 재무]}{통제관 \textbf{P1 PM} / 데이터 \textbf{박세연}(게이트 2 판정)·전환 용역사 PM / 인터페이스 수행 인터페이스 PL(후임)·API GW 벤더 담당(원격 대기) / 매장 배포 \textbf{매장운영팀장}·지역매니저 4명(강윤지 후임 포함) / 검증 P1 PMO 품질 담당 / 롤백 \textbf{이재현} / 운영 인수 \textbf{박준영}·서비스데스크 리더 / 규제 최영주(원격, PG·개인정보) / 기록 P1 PMO 기준선 담당 / 재무 재경팀장(원격, 결제 대사)} \\ \hline
\B{\rd{14mm}}{}판정 권한 & \Bg{[게이트 1 실명(위임 근거) / 게이트 2 / 게이트 3·오픈 / PoNR 이후 롤백 / 판정자 부재 시 대체]}{게이트 1 P1 PM(정미란 위임장 02-19) / 게이트 2 박세연 / 게이트 3·오픈 \textbf{정미란}(02-27 11:00 현장 도착) / PoNR(01:40) 이후 손실 있는 롤백 정미란(부재 시 김선호 — 헌장 §14 중단 조건과 같은 층) / 판정자 부재 시 대체: 게이트 2 전환 용역사 PM + P1 PM 합의, 게이트 3 박세연} \\ \hline
\B{\rd{9mm}}{}인력 · 교대 & \Bg{[총원, 갈래별 교대, 통제관 수면·대행]}{82명 전원 + 서비스데스크 11 + 지역매니저 4 = 97명. 교대: 데이터 갈래 2교대(20:00~08:00 / 08:00~20:00), 통제관 P1 PM 연속(수면 03:00~06:00, 대행 P1 PMO장)} \\ \hline
\B{\rd{8mm}}{}연락망 & \Bg{[유선·채널·매장 긴급 회선·벤더 원격 회선 — 메신저 하나뿐이면 안 된다]}{컨트롤룸 유선 2회선·컨트롤룸 채널·매장 긴급 회선(서비스데스크 11명 3교대)·벤더 원격 회선 3(API GW·MDM·CSP)} \\ \hline
\B{\rd{12mm}}{}리허설 & \Bg{[1회차 일자·소요·결과 / 2회차 / 3회차 → 계획 대비 여유]}{1회차 11-23 \textbf{41h 20m}(7.4TB 전체, 창 초과, 검증 3종 미통과) / 2회차 01-08 \textbf{33h 10m}(5.1TB, 6종 통과, 표본 3건 보정) / 3회차 01-22 \textbf{31h 40m}(6종 통과, 롤백 리허설 7h 50m 병행) → 계획 34h 대비 여유 2h 20m} \\ \hline
\B{\rd{10mm}}{}입력 & \Bg{[게이트 go 기록, 시험 종료 보고, 리허설 검증 보고, 연결 시험, 배포 리허설, 매장 통보]}{G3 go(02-19), UAT 종료 보고(02-18, 심각도 1·2 미해결 0), 리허설 3회차 검증 6종 보고, 인터페이스 47본 연결 시험(02-10), 펌웨어 배포 리허설(파일럿 12개점 + 구형 2점 100\%, 02-09), 매장 통보(02-12 1회 통합 — R-14 대응)} \\ \hline
\end{tabularx}}

% ------------------------------------------------------------------ 2 (가로) 시각표
\landscapeon
\secbar{2. 시각표 — T$-$24h ~ T+24h \B{}{(계획 / 실행 기록)}}
\guide{행마다 계획 시각(T 기준·실제 시각), 갈래(D 데이터 · I 인터페이스 · S 매장 · C 통제), 작업, 담당(실명), \textbf{확인 = 관측 가능한 값}(건수·해시·응답 코드·성공률 — “정상 확인”은 확인이 아니다), 롤백 판단점 ★(이 행에서 되돌릴 수 있는가·비용), 실행 기록(실제 시각·결과·편차). 리허설 실측을 계획 시각의 근거로 쓴다. 게이트와 PoNR은 별도 행으로 굵게. 병렬 갈래는 같은 시각의 행을 갈래로 구분한다. 실행 기록 열은 본 컷오버에서 채운다 — 이 열이 다음 런북의 근거다.}
{\scriptsize\setlength{\tabcolsep}{2pt}
\begin{longtable}{|C{7mm}|L{13mm}|L{19mm}|C{7mm}|L{63mm}|L{25mm}|L{45mm}|L{30mm}|L{44mm}|}\hline
\hd{\#} & \hd{T} & \hd{계획 시각} & \hd{갈래} & \hd{작업} & \hd{담당} & \hd{확인 (관측값 · 기준)} & \hd{★ 롤백 판단점 (비용)} & \hd{실행 기록 (실제 · 결과)} \\ \hline \endhead
\ifwbblank
\rd{8mm}1 & \g{T$-$24h} & & \g{C} & \g{[최종 상태 점검 회의 — 게이트 이후 변화 확인, 런북 배포]} & & \g{[신규 심각도 1·2 = 0, 변경 동결 확인]} & & \\ \hline
\rd{8mm}2 & \g{T$-$22h} & & \g{D} & \g{[구 시스템 최종 전 백업 1차]} & & \g{[스냅샷 해시]} & & \\ \hline
\rd{8mm}3 & \g{T$-$12h} & & \g{C} & \g{[컨트롤룸 개설 · 연락망 · 롤백 패키지 해시]} & & & & \\ \hline
\rd{8mm}4 & \g{T$-$10h} & & \g{S} & \g{[배포 서버 · 배포 큐 적재]} & & \g{[큐 n/n, 패키지 해시]} & & \\ \hline
\rd{8mm}5 & \g{T$-$6h} & & \g{I} & \g{[외부 상대(3PL·PG·배달·앱) 전환 시각 재확인]} & & \g{[회신 n/n]} & & \\ \hline
\rd{8mm}6 & \g{T$-$2h} & & \g{C} & \g{[인력 출석 · 라이선스 유효 · 매장 안내]} & & \g{[출석 n/n, 키 응답]} & & \\ \hline
\rd{8mm}\textbf{7} & \g{T$-$0.5h} & & \g{C} & \textbf{게이트 1 — 착수 go/no-go} \g{(항목 3)} & & \g{[조건 n개 전부 충족]} & \g{★ 비용 0} & \g{[go 시각 · 서명]} \\ \hline
\rd{8mm}8 & \g{T0} & & \g{S} & \g{[매장 거래 정지 → 오프라인 모드]} & & \g{[전환 응답 n/n]} & & \\ \hline
\rd{8mm}9 & & & \g{D} & \g{[구 시스템 최종 마감 배치]} & & \g{[종료 코드 · 마감 건수]} & & \\ \hline
\rd{8mm}10 & & & \g{D} & \g{[최종 백업 2차 + 전환 대상 추출]} & & \g{[행 수 = 리허설 ±\%]} & & \\ \hline
\rd{8mm}11 & & & \g{D} & \g{[마스터 전환]} & & & & \\ \hline
\rd{8mm}12 & & & \g{I} & \g{[인터페이스 구 연결 차단]} & & \g{[차단 n/n]} & & \\ \hline
\rd{8mm}13 & & & \g{D} & \g{[검증 6종 1차(마스터)]} & & \g{[6종 통과 — 항목 6]} & & \\ \hline
\rd{8mm}\textbf{14} & & & \g{D} & \textbf{게이트 2 — 데이터 정본 전환 go/no-go} & & \g{[6종 통과, 재실행 ≤ n회]} & \g{★ 비용 = 시간 (PoNR 전)} & \\ \hline
\rd{8mm}\textbf{15} & & & \g{D} & \textbf{Point of no return — 신 원장 활성화} \g{(항목 4)} & & \g{[활성 플래그 · 구 원장 읽기 전용]} & \g{★★ 이후 롤백 = 손실} & \\ \hline
\rd{8mm}16 & & & \g{D} & \g{[이력 적재 시작 — 리허설 소요 n h]} & & \g{[진행률 보고 주기, 미완 시 RB-n]} & & \\ \hline
\rd{8mm}17 & & & \g{I} & \g{[첫 실운영 배치(폴백 등)]} & & & & \\ \hline
\rd{8mm}18 & & & \g{C} & \g{[교대 · 중간 보고]} & & & & \\ \hline
\rd{8mm}19 & & & \g{D} & \g{[이력 적재 완료 예정]} & & \g{[적재 행 수 = 추출 행 수]} & & \\ \hline
\rd{8mm}20 & & & \g{D} & \g{[검증 6종 2차(이력)]} & & & & \\ \hline
\rd{8mm}21 & & & \g{I} & \g{[인터페이스 신 연결 전환 · 연결 시험]} & & \g{[n/n 응답 정상]} & & \\ \hline
\rd{8mm}22 & & & \g{I} & \g{[성능 스모크]} & & \g{[p95 < n ms]} & & \\ \hline
\rd{8mm}\textbf{23} & & & \g{C} & \textbf{게이트 3 — 매장 배포 go/no-go} & & & \g{★ 비용 = 부분 롤백 n억·n일} & \\ \hline
\rd{8mm}24 & & & \g{S} & \g{[매장 배포 시작 — 병렬 n, 매장당 n분]} & & \g{[진행률 시간당 보고]} & & \\ \hline
\rd{8mm}25 & & & \g{C} & \textbf{기술적 롤백 최종 시한} \g{(오픈 $-$ 롤백 소요)} & & \g{[트리거 관측값 확인]} & \g{★ 마지막 전체 롤백 시각} & \\ \hline
\rd{8mm}26 & & & \g{S} & \g{[배포 결과 집계]} & & \g{[실패 ≤ n점 (RB-n)]} & \g{★ 부분 롤백} & \\ \hline
\rd{8mm}27 & & & \g{S} & \g{[실패 재배포]} & & & & \\ \hline
\rd{8mm}28 & & & \g{S} & \g{[파일럿 스모크 거래]} & & \g{[성공률 ≥ n\%]} & & \\ \hline
\rd{8mm}29 & & & \g{D} & \g{[재고 원장 대사]} & & \g{[차이 0]} & & \\ \hline
\rd{8mm}30 & & & \g{C} & \g{[최종 상태 확인 — 대시보드 · 데스크 배치]} & & \g{[경보 0]} & & \\ \hline
\rd{8mm}\textbf{31} & & & \g{C} & \textbf{오픈 판정} & & & \g{★ 오픈 연기} & \g{[go 시각 · n인 서명]} \\ \hline
\rd{8mm}\textbf{32} & \g{T+nh} & & \g{C} & \textbf{오픈} & & \g{[온라인 전환 응답 n/n]} & & \\ \hline
\rd{8mm}33 & \g{D+1} & & \g{C} & \g{[하이퍼케어 1일차 — 시간대별 성공률 · 콜 · 첫 일마감]} & & & & \\ \hline
\rd{8mm}34 & \g{D+1} & & \g{C} & \textbf{D+1 판정 회의 → ELS 착수 인계} \g{(⑰)} & \g{[→ 인수자]} & \g{[ELS 착수 조건]} & & \g{[인수 서명 시각]} \\ \hline
\rd{8mm} & & & & & & & & \\ \hline
\rd{8mm} & & & & & & & & \\ \hline
\else
1 & T$-$24h & 02-25(목) 20:00 & C & 최종 상태 점검 회의: G3 이후 변화(신규 심각도 1·2 결함, 인프라 변경, 인력) 확인, 런북 v1.1 배포 & P1 PM & 신규 심각도 1·2 = 0, 변경 동결(CAB 잠금 02-20~03-12) 확인 & — & 20:00~20:50, 변화 없음. D$-$2 라이선스 사건 보고(항목 9) \\ \hline
2 & T$-$22h & 22:00 & D & 구 시스템(OD-POS 2.1·SAP) 최종 전 백업 1차(원본 7.4TB 스냅샷) & 박세연·인프라팀 & 스냅샷 해시 일치, 복원 시험 이력(01-22) 참조 & — & 22:00~01:10 완료 \\ \hline
3 & T$-$12h & 02-26(금) 08:00 & C & 컨트롤룸 개설, 연락망 점검(유선·채널·매장 회선), 롤백 패키지 해시 확인 & P1 PM·이재현 & 롤백 패키지 3종(구 POS 이미지·SAP IF 설정·DNS) 해시 = 01-22 리허설본 & — & 08:00 개설, 해시 일치 \\ \hline
4 & T$-$10h & 10:00 & S & 615개 매장 배포 서버 3대 상태·배포 큐 적재(615 패키지, 구형 기종 35점 별도 패키지 CR-14) & 매장운영팀장·수행 POS팀 & 큐 615/615, 패키지 해시 2종 일치 & — & 10:20 완료 \\ \hline
5 & T$-$6h & 14:00 & I & 3PL 대성로지스 통보(폴백 배치 21:00·03:00 → 컷오버 주말 03:00 1회차 신 원장 기준), 배달3사·PG 7본·앱 벤더 전환 시각 재확인 & 수행 인터페이스 PL·한동석 & 회신 확인 3PL/배달3사/PG/앱 4/4 & — & 14:00~15:30 회신 완료 \\ \hline
6 & T$-$4h & 16:00 & D & 가맹 발주 포털 마감 20:00 유지 공지(22:00 파라미터 미적용, CR-09 조건 ①), 당일 발주분 구 시스템 처리 & 오정근 대리(매장운영팀장) & 포털 공지 게시, 20:00 마감 설정 확인 & — & 16:00 게시 \\ \hline
7 & T$-$2h & 18:00 & C & 인력 출석 점검, 임시 라이선스 유효 확인(만료 03-26), 매장 오프라인 모드 안내 재발송(로컬 캐시 판매, 새벽 동기화 [추가 설정]) & P1 PM·강현도 & 출석 97/97, 라이선스 키 유효 응답 200 & — & 18:00 출석 96/97(지역매니저 1명 19:10 도착), 키 유효 \\ \hline
\textbf{8} & \textbf{T$-$0.5h} & \textbf{19:30} & \textbf{C} & \textbf{게이트 1 — 착수 go/no-go} (항목 3) & \textbf{P1 PM} & 조건 6개 전부 충족 & ★ (비용 0 — 착수 안 하면 됨) & \textbf{19:35 go}, 서명 P1 PM·이재현·박준영 \\ \hline
9 & \textbf{T0} & 20:00 & S & 615개 매장 OD-POS 거래 정지 → 오프라인 모드 전환(22:00 마감 매장은 로컬 캐시 판매), 직영 227 마감 순차 & 매장운영팀장 & 오프라인 모드 전환 응답 615/615(20:00~20:40) & — & 20:00~20:45, 615/615 \\ \hline
10 & T+0.5h & 20:30 & D & 구 시스템 최종 마감 배치(OD-POS 일마감·가맹 발주 마감 처리·SAP 일마감) & 박세연·SAP팀 & 마감 배치 3종 정상 종료 코드, 마감 건수 기록 & — & 20:30~21:25 \\ \hline
11 & T+1.5h & 21:30 & D & 최종 백업 2차(마감 후 상태) + 전환 대상 추출(5년 필터, 5.1TB) 시작 & 전환 용역사 PM & 스냅샷 해시, 추출 행 수 = 리허설 3회차 ±0.5\% & — & 21:30~22:05 완료, 행 수 +0.2\%(2월 거래 증가분) \\ \hline
12 & T+2h & 22:00 & D & 마스터 6종 전환(상품·고객·매장·거래처·가격·재고, 2,184,310행) & 전환 용역사 PM & 적재 완료 로그, 행 수 일치 & — & 22:00~00:20(계획 00:30) \\ \hline
13 & T+2h & 22:00 & I & 인터페이스 47본 구 연결 차단(폐기 2본 차단 목록 포함), 신 연결 설정 대기 상태 & 수행 인터페이스 PL & 차단 47/47, 폐기 2본 차단 확인 & — & 22:00~22:40 \\ \hline
14 & T+4.5h & 02-27(토) 00:30 & D & \textbf{검증 6종 1차(마스터)} — 건수·금액·키·참조·매핑·표본 2,000행 & P1 PMO 품질·전환 용역사 & 6종 통과(항목 6) & — & 00:30~01:00: 5종 통과, \textbf{표본 대조 1건 불일치}(가격 마스터 VAT 포함가 반올림) \\ \hline
\textbf{15} & \textbf{T+5h} & \textbf{01:00} & \textbf{D} & \textbf{게이트 2 — 데이터 정본 전환 go/no-go} (항목 3) & \textbf{박세연} & 6종 통과, 재실행 ≤ 2회 & ★ (비용 = 시간. PoNR 전, 손실 0) & 01:00 보류 → 변환 규칙 수정·가격 마스터 재실행 01:12~01:35 → 표본 0건 → \textbf{01:38 go}(+38분) \\ \hline
\textbf{16} & \textbf{T+5h40m} & \textbf{01:40} & \textbf{D} & \textbf{Point of no return — 신 재고·주문 원장 활성화}(정본 전환), 이후 거래·배치는 신 원장에만 기록 & 박세연·P1 PM & 원장 활성 플래그, 구 원장 읽기 전용 전환 확인 & ★★ 이후 롤백 = 손실 있음(항목 4), 결정권 정미란 & \textbf{01:40 활성화}(계획과 동일 — 게이트 2 지연분은 여유에서 흡수) \\ \hline
17 & T+5h40m & 01:40 & D & 거래 이력 5년 적재 시작(3.21억 행, 5.1TB, 리허설 6h 50m) & 전환 용역사 PM & 진행률 30분 단위 보고, 11:00 미완 시 RB-2 & — & 01:40 시작 \\ \hline
18 & T+7h & 03:00 & I & 3PL 폴백 배치 1회차 — 신 원장 기준 출고 지시 생성·전송(첫 실운영) & 수행 인터페이스 PL·이천 상황실 & 배치 종료 40분 이내, 대성 수신 확인 & — & 03:00~03:31, 수신 확인 03:40 \\ \hline
19 & T+8h & 04:00 & C & 교대 1(통제관 대행 P1 PMO장 03:00~06:00), 중간 보고 1(스폰서·프로그램 PMO장 문자) & P1 PMO장 & — & — & 04:00 보고 “게이트 2 +38, PoNR 정시” \\ \hline
20 & T+12h & 08:00 & C & 데이터 갈래 2교대 인수, 통제관 복귀 & P1 PM & 인수 체크리스트 & — & 08:00 \\ \hline
21 & T+13h & 09:00 & D & 이력 적재 완료 예정 & 전환 용역사 PM & 적재 행 수 = 추출 행 수 & — & \textbf{08:35 완료}(6h 55m, 리허설 +5분) \\ \hline
22 & T+13h & 09:00 & D & \textbf{검증 6종 2차(이력)} & P1 PMO 품질·전환 용역사·재경팀장(금액) & 6종 통과, 매출 원장 5년 합계 일치 & — & 08:40~10:50 전부 통과(표본 0건) \\ \hline
23 & T+15h & 11:00 & I & 인터페이스 47본 신 연결 전환·연결 시험(SAP 애드온 8·PG 7·배달3사 3·앱 1·멤버십·BI·기타) & 수행 인터페이스 PL·API GW 벤더 & 47/47 응답 정상, PG 7본 승인 시험 거래 1건씩 & — & 11:00~12:15, 47/47(PG 1본 타임아웃 1회 → 재시도 정상) \\ \hline
24 & T+15h & 11:00 & I & 성능 스모크(재고 API 응답 p95 < 800ms, 주문 허브 처리 100건/분) & 수행 시험팀 & 부하시험 기준(01월) 대비 & — & 11:20~12:00 p95 610ms \\ \hline
\textbf{25} & \textbf{T+16h} & \textbf{12:00} & \textbf{C} & \textbf{게이트 3 — 매장 배포 go/no-go} (항목 3) & \textbf{정미란} & 검증 2차 6종·인터페이스 47/47·스모크 통과 & ★ (비용 = 배포 후 매장 부분 롤백 2.1억·6영업일) & \textbf{12:30 go}(+30분, 정미란 11:00 도착·검증 보고 열람), 서명 정미란·박세연·P1 PM \\ \hline
26 & T+17h & 13:00 & S & POS 펌웨어·SW 일괄 배포 시작(배포 서버 3대, 8병렬, 매장당 약 15분 → 615점 7시간), 구형 35점 별도 패키지 & 매장운영팀장·수행 POS팀 & 배포 진행률 시간당 보고, 실패 목록 실시간 & — & 13:00 시작 \\ \hline
27 & T+17h & 13:00 & D & CR-08 임시 연동 2본(파일럿 재고 동기화) 신 원장 연결 유지 확인 — 폐기는 2차 후 30일(04-27, CCB 01-21 갱신) & 박세연 & 연동 2본 응답 정상 & — & 13:10 확인 \\ \hline
28 & T+20h & 16:00 & C & \textbf{기술적 롤백 최종 시한}(오픈 06:00 $-$ 14h) — 이후 롤백 불가, 오픈 연기만 가능 & P1 PM·정미란 & 트리거 RB-1~3 관측값 확인 & ★ 마지막 전체 롤백 시각 & 16:00 관측: 배포 412/615 진행, 실패 4, 트리거 해당 없음 → 롤백 창 종료 \\ \hline
29 & T+24h & 20:00 & S & 배포 1차 결과 집계 & 매장운영팀장 & 실패 ≤ 31점(5\%)·직영 ≤ 10점(RB-3) & ★ 부분 롤백(매장 갈래) & \textbf{20:10 완료, 실패 9점}(직영 2·가맹 7, 구형 기종 5 포함) — RB-3 미발동 \\ \hline
30 & T+24h30m & 20:30 & S & 실패 9점 재배포 & 지역매니저 4명·수행 POS팀 & 재배포 성공 & — & 20:30~22:40 \textbf{8점 성공, 1점(가맹 구형 2011년 기종) 수기 모드 → 03-02 방문 교체(CR-14)} \\ \hline
31 & T+27h & 23:00 & S & 파일럿 12개점 + 지역매니저 스모크 거래(실거래 1건·취소 1건·반품 1건, 멤버십 적립·PG 승인 포함) & 지역매니저 4명·강윤지 후임 & 성공률 ≥ 99.0\%, PG 7본 0 실패(RB-4) & — & 23:00~01:30, 36건 중 36 성공(100\%) \\ \hline
32 & T+29h & 01:00 & D & 재고 원장 대사(마스터 재고 vs 매장 오프라인 캐시 판매분 동기화 615점) & 박세연·수행 재고팀 & 동기화 615/615, 대사 차이 0 & — & 01:00~02:40, 615/615, 차이 0 \\ \hline
33 & T+31h & 03:00 & I & 3PL 폴백 배치 2회차(신 원장, 일요일 출고 없음 — 빈 배치 정상 종료 확인) & 수행 인터페이스 PL & 정상 종료 코드 & — & 03:00~03:08 \\ \hline
34 & T+32h & 04:00 & C & 최종 상태 확인: 모니터링 대시보드(APM) 임계값 0 초과, 서비스데스크 3교대 배치, 매장 안내 문자 예약 & 박준영·P1 PM & 대시보드 경보 0, 데스크 11/11 & — & 04:00 경보 0 \\ \hline
\textbf{35} & \textbf{T+32h30m} & \textbf{04:30} & \textbf{C} & \textbf{오픈 판정} — RB-4 관측, 게이트 3 이후 변화 없음 & \textbf{정미란}(현장)·박준영·이재현·P1 PM & 스모크 100\%, PG 0 실패, 배포 614/615(1점 수기 승인) & ★ 오픈 연기 & \textbf{04:30 go}, 4인 서명 \\ \hline
36 & \textbf{T+34h} & \textbf{06:00} & \textbf{C} & \textbf{오픈} — 615개 매장 온라인 모드 전환, 직영 조기 매장 07:00 영업 & 매장운영팀장 & 온라인 전환 응답 614/615 & — & \textbf{06:00 오픈}, 06:00~06:30 전환 614/615 \\ \hline
37 & D+1 & 02-28(일) 06:00~22:00 & C & 하이퍼케어 1일차: 시간대별 거래 성공률(목표 99.0\%)·재고 API p95·PG 승인율·서비스데스크 콜, 12:00 중간 보고, 22:00 신 플랫폼 첫 일마감 & 박준영·이재현 & 거래 성공률 시간대별 ≥ 99.0\% & — & 성공률 최저 99.3\%(11시대), 콜 214건(주문 오류·조작 문의 71\%), 22:00 일마감 정상 \\ \hline
38 & D+1 & 03-01(월) 07:00 & C & \textbf{D+1 영업일 판정 회의 → ELS 착수 인계}(⑰), 컨트롤룸 → 워룸 축소, 임시 라이선스 정식 결재 추적(03-05) & P1 PM → \textbf{박준영} & ELS 착수 조건: 일마감 정상·03:00 배치 정상·경보 0 & — & 07:00 회의, \textbf{박준영 인수 서명 07:20}, ELS 1주차 착수 \\ \hline
\fi
\end{longtable}}
\B{\small\g{계획 대비 편차(게이트·적재·배포·오픈, 분 단위) / 여유 소진: }\hrulefill\par}{\small \textbf{계획 대비 편차}: 게이트 2 +38분, 게이트 3 +30분, 이력 적재 $-$25분, 배포 완료 +10분, 오픈 정시. 총 여유 2h 20m 중 소진 약 1h 10m.\par}

% ------------------------------------------------------------------ 3 (가로) go/no-go 게이트
\secbar[70mm]{3. go/no-go 게이트 — 조건 · 판정자 · 시각}
\guide{게이트는 되돌리는 비용이 다른 세 지점에 선다 — 착수(비용 0), 데이터 정본 전환 직전(비용 = 시간), 매장 배포 직전(비용 = 돈과 매장) + 오픈 판정. 조건은 전부 충족형·관측 가능한 값으로 적고 “재실행 허용 횟수”를 붙인다. 판정자는 실명·위임·대체자. 미충족 시 선택지(대기·재실행·롤백·부분 진행)와 각각의 \textbf{시간 한도}를 미리 적는다 — 새벽 3시의 “PM 판단”은 판단이 아니라 피로다. 판정 기록(시각·판정·서명) 칸을 둔다.}
{\footnotesize\setlength{\tabcolsep}{2.5pt}
\begin{tabularx}{\linewidth}{|L{24mm}|L{26mm}|Y|C{13mm}|L{30mm}|L{52mm}|L{40mm}|}\hline
\hd{게이트} & \hd{계획 / 실제} & \hd{조건 (전부 충족)} & \hd{재실행} & \hd{판정자 (대체)} & \hd{미충족 시 선택지 (시간 한도)} & \hd{판정 기록} \\ \hline
\ifwbblank
\rd{13mm}\textbf{1 착수} & \g{[hh:mm / ]} & \g{① 인력 출석 ② 백업 해시 ③ 라이선스·결재 상태 ④ 매장 통보 ⑤ 이전 게이트 이후 변화 없음 ⑥ 외부(기상·PG 점검)} & \g{—} & \g{[실명(위임)]} & \g{[착수 연기 최대 hh:mm → 이후 다음 창]} & \g{[go 시각 · 서명]} \\ \hline
\rd{13mm}\textbf{2 데이터 정본 전환} & & \g{① 마스터 검증 n종 통과 ② 재실행 n회 이내 ③ 구 원장 읽기 전용 준비} & \g{[n회]} & & \g{[재실행 초과 → RB-1 전체 롤백 / 지연 시 후속 배치 이동]} & \\ \hline
\rd{13mm}\textbf{3 매장 배포} & & \g{① 이력 검증 통과 ② 인터페이스 n/n ③ 성능 스모크 ④ 패키지 해시 ⑤ RB-2 미발동} & & & \g{[재시험 후 재판정 hh:mm까지, 초과 시 롤백 또는 연기]} & \\ \hline
\rd{13mm}오픈 판정 & & \g{① 스모크 ≥ n\% ② 배포 ≥ n\% ③ 대사 차이 0 ④ 경보 0} & & & \g{[오픈 연기 최대 hh:mm]} & \\ \hline
\else
\textbf{1 착수} & 02-26 19:30 / 19:35 & ① 인력 97명 출석(핵심 역할 12명 전원) ② 백업 1차 해시 일치 ③ \textbf{임시 라이선스 유효(만료 03-26 > 정식 03-05)} ④ 매장 통보·오프라인 모드 안내 완료 ⑤ G3 이후 신규 심각도 1·2 결함 0, CAB 잠금 유지 ⑥ 외부: 기상 특보·PG사 점검 공지 없음 & — & P1 PM(정미란 위임) & 미충족 시 착수 연기 최대 02-26 22:00(이후 창 34h 확보 불가 → 대안 2 03-12) & go 19:35, 서명 3인. ①은 19:10 충족 \\ \hline
\textbf{2 데이터 정본 전환} & 02-27 01:00 / \textbf{01:38} & ① 마스터 6종 검증 6종 통과 ② 재실행 2회 이내 ③ 구 원장 읽기 전용 전환 준비 & 2회 & \textbf{박세연}(대체: 용역사 PM + P1 PM) & 재실행 2회 초과 → RB-1 전체 롤백(PoNR 전, 손실 0) / 03:00 이후 통과 시 03:00 폴백 배치를 04:30로 이동(3PL 통보) & 표본 대조 1건 불일치 → 재실행 1회 → go 01:38 \\ \hline
\textbf{3 매장 배포} & 02-27 12:00 / \textbf{12:30} & ① 이력 검증 6종 통과 ② 인터페이스 47/47 ③ 성능 스모크 통과 ④ 배포 패키지 해시 확인 ⑤ RB-2 미발동 & — & \textbf{정미란}(대체: 박세연) & 미충족 시 13:00~16:00 내 재시험 후 재판정, 16:00 초과 시 전체 롤백 또는 오픈 연기 결정 & go 12:30, 서명 정미란·박세연·P1 PM \\ \hline
오픈 판정 & 02-28 04:30 / 04:30 & ① 스모크 ≥ 99.0\%·PG 0 실패 ② 배포 ≥ 95\%(수기 모드 매장 목록 승인) ③ 재고 대사 차이 0 ④ 경보 0 & — & 정미란·박준영·이재현·P1 PM & 오픈 연기(06:00 → 최대 10:00, 매장 수기 운영) & go 04:30 \\ \hline
\fi
\end{tabularx}}

% ------------------------------------------------------------------ 4 Point of no return
\secbar[70mm]{4. Point of no return — 시각과 근거}
\guide{두 시각을 분리한다. \textbf{PoNR} = 신 원장이 ‘정본’이 되어 이후 거래가 신 원장에만 기록되기 시작하는 시각. \textbf{기술적 롤백 최종 시한} = 오픈 $-$ (구 시스템 복원 + 검증 + 매장 캐시 재동기화). 둘 사이의 롤백은 “가능하지만 손실 있음”이며 결정권자를 한 층 올린다. PoNR을 “배포 시작”으로 적지 않는다 — 매장 배포는 되돌릴 수 있고 원장 분기는 되돌리기 어렵다.}
{\footnotesize
\begin{tabularx}{\linewidth}{|L{30mm}|Y|}\hline
\hd{항목} & \hd{내용} \\ \hline
\B{\rd{8mm}}{}PoNR 시각 & \Bg{[YYYY-MM-DD hh:mm (T+n) — 계획 / 실제]}{\textbf{2027-02-27(토) 01:40} (T+5h40m) — 계획·실제 동일} \\ \hline
\B{\rd{12mm}}{}근거 & \Bg{[무엇이 정본으로 바뀌는가 — 이 시각부터 신 원장에만 기록되는 것 ①②③]}{게이트 2 통과 직후 신 재고·주문 원장을 ‘정본’으로 활성화한다. 이 시각부터 ① 이력 적재가 신 원장에 기록되고 ② 03:00 3PL 폴백 배치 1회차가 신 원장 기준으로 출고 지시를 생성해 대성로지스에 전송되며 ③ 매장 오프라인 캐시 판매분이 신 원장으로 동기화된다. 즉 01:40 이후 구 시스템으로 되돌리면 신 원장에만 있는 거래·출고 지시가 생긴다} \\ \hline
\B{\rd{10mm}}{}이후 롤백이 잃는 것 & \Bg{[대사 대상 건수(리허설 추정)·영업일, 외부 협의, 데이터 손실 유무]}{03:00 배치 출고 지시(리허설 추정 1,200건 이상 — 토요일 새벽 출고분)와 매장 캐시 판매분의 구 원장 수작업 대사 2영업일, 대성로지스 출고 취소 협의. 데이터 손실은 없으나 이중 원장 정리 비용 발생 [추가 설정]} \\ \hline
\B{\rd{9mm}}{}기술적 롤백 최종 시한 & \Bg{[오픈 hh:mm $-$ n시간 (복원 n h + 검증 n h + 재동기화 n h, 롤백 리허설 실측)]}{\textbf{02-27(토) 16:00} = 오픈 06:00 $-$ 14시간(구 시스템 복원 8h + 검증 2h + 매장 캐시 재동기화 4h, 롤백 리허설 01-22 실측 7h 50m + 여유)} \\ \hline
\B{\rd{9mm}}{}롤백 결정권자 & \Bg{[PoNR 전 실명 / PoNR ~ 기술적 시한 실명(부재 시) / 시한 이후 — 롤백 불가, 연기만]}{01:40 이전 P1 PM(게이트 1·2 판정 권한 범위) / 01:40 ~ 16:00 \textbf{정미란}(부재 시 김선호) / 16:00 이후 롤백 불가, 오픈 연기만(정미란)} \\ \hline
\B{\rd{7mm}}{}리허설 검증 & \Bg{[롤백 리허설 일자·시나리오·소요·대사 대상]}{롤백 리허설(01-22, 3회차 병행): PoNR 후 롤백 시나리오 7h 50m, 대사 대상 1,140건 추정} \\ \hline
\end{tabularx}}

% ------------------------------------------------------------------ 5 롤백 트리거
\secbar[100mm]{5. 롤백 트리거 — 임계값 · 판정자 · 발동 시 조치}
\guide{트리거 4개 안팎. 관측 지표·\textbf{숫자 임계값}(리허설 실측 + 허용범위에서 온다)·관측 시각과 시각표 행·판정자(그 시각에 컨트롤룸에 있는 사람)·발동 시 조치(전체 롤백 / 부분 롤백 / 오픈 연기)·복구 소요. 부분 롤백이 가능한 트리거와 불가능한 트리거를 구분한다. 발동하지 않은 트리거도 “관측값 / 임계값”을 적는다 — 다음 컷오버의 임계값 근거다. 롤백 후 다음 창을 적는다.}
{\footnotesize\setlength{\tabcolsep}{2.5pt}
\begin{tabularx}{\linewidth}{|C{10mm}|L{30mm}|L{36mm}|L{26mm}|L{28mm}|Y|L{26mm}|L{36mm}|}\hline
\hd{ID} & \hd{지표} & \hd{임계값} & \hd{관측 시각 · 행} & \hd{판정자} & \hd{발동 시 조치 (전체 / 부분 / 연기)} & \hd{복구 소요} & \hd{본 컷오버 관측값 / 판정} \\ \hline
\ifwbblank
\rd{11mm}RB-1 & \g{[마스터 검증 미통과]} & \g{[재실행 n회 후 n종 이상]} & \g{[게이트 2, 행 n]} & & \g{[전체 롤백 — PoNR 전, 손실 0]} & & \g{[관측값 / 임계값 → 발동 여부]} \\ \hline
\rd{11mm}RB-2 & \g{[이력 적재 완료 시각]} & \g{[리허설 + n h 미완]} & & \g{[PoNR 후 → 상위]} & \g{[손실 있는 롤백 vs 오픈 연기]} & & \\ \hline
\rd{11mm}RB-3 & \g{[배포 실패 매장 수]} & \g{[n점(n\%) 또는 직영 n점]} & & & \g{[매장 갈래 부분 롤백 — 데이터 유지]} & & \\ \hline
\rd{11mm}RB-4 & \g{[스모크 성공률 · PG]} & \g{[n\% 미만 또는 n본 실패]} & \g{[오픈 판정]} & & \g{[오픈 연기 → 부분 운영]} & & \\ \hline
\rd{11mm} & & & & & & & \\ \hline
\else
RB-1 & 마스터 검증 6종 중 미통과 항목 & 재실행 2회 후 1종 이상 미통과 & 게이트 2(01:00), 행 14~15 & 박세연 & \textbf{전체 롤백}(PoNR 전 — 구 시스템 재기동 2h, 손실 0), 다음 창 대안 2(03-12~14) 준비 & 2h & 표본 대조 1건 → 재실행 1회 후 0건 / \textbf{미발동} \\ \hline
RB-2 & 이력 적재 완료 시각 & 02-27 11:00 미완(리허설 6h 50m + 4h) & 행 17~21 & P1 PM → 정미란(PoNR 후) & 손실 있는 롤백(대사 2영업일) vs 오픈 연기(최대 02-28 10:00) — 정미란 판단 & 롤백 8h+ / 연기 4h & 08:35 완료 / \textbf{미발동} \\ \hline
RB-3 & 펌웨어 배포 실패 매장 수 & 31점(5\%) 초과 또는 직영 10점 초과 & 행 29(20:00) & 오정근(부재 시 매장운영팀장)·박세연 & \textbf{매장 갈래 부분 롤백}(실패 매장 구 POS 이미지 복원, 데이터는 신 원장 유지), 오픈 연기 검토 & 매장당 30분 원격 / 방문 시 6영업일·2.1억 & 실패 9점(직영 2) / \textbf{미발동} \\ \hline
RB-4 & 스모크 거래 성공률·PG 승인 & 99.0\% 미만 또는 PG 7본 중 1본 이상 실패 & 오픈 판정(04:30), 행 31·35 & 정미란 & 오픈 연기(최대 10:00, 매장 수기) → 원인 미해결 시 전체 롤백 불가(16:00 경과) → 신 플랫폼 부분 운영(PG 대체 경로) & 연기 4h & 36/36(100\%), PG 0 실패 / \textbf{미발동} \\ \hline
\fi
\end{tabularx}}
\B{\small\g{발동 시 다음 창(EX 대안·합류 버퍼) / 발동 0이어도 관측값 기록 이유: }\hrulefill\par}{\small 트리거 발동 0. 발동 시 다음 창 = EX-01 대안 2(03-12~14, 합류 버퍼 15영업일) 재소환. 관측값은 다음 컷오버(2차 03-27, 향후 R5·P2)의 임계값 근거로 기록.\par}

% ------------------------------------------------------------------ 6 (가로) 데이터 검증 6종
\secbar[110mm]{6. 데이터 검증 체크리스트 6종 — 분모와 통과 기준 \B{}{(1,842테이블 / 원본 7.4TB → CR-11 5년 필터 후 5.1TB; 마스터 6종 2,184,310행, 이력 3.21억 행)}}
\guide{건수·금액 합계·키 유일성·참조 정합성·코드 매핑 커버리지·표본 필드 대조 — 마스터(게이트 2)와 이력(게이트 3) 두 번 적용한다. \textbf{분모를 명시}한다(전환 대상 테이블 수·필터 후 볼륨). 통과 기준은 0이 원칙, 표본 대조는 재실행 허용 횟수를 붙인다. 리허설 1·2·3회차의 \textbf{실패 항목과 조치}를 옆에 적는다 — 3회차가 통과했다는 사실보다 1회차에서 무엇이 실패했는지가 본 컷오버의 위험을 말해 준다. 표본은 무작위로 뽑는다.}
{\footnotesize\setlength{\tabcolsep}{2pt}
\begin{tabularx}{\linewidth}{|L{22mm}|L{30mm}|L{28mm}|L{18mm}|L{34mm}|L{24mm}|L{16mm}|L{34mm}|L{22mm}|Y|}\hline
\hd{검증} & \hd{대상 · 분모} & \hd{산식} & \hd{통과 기준} & \hd{리허설 1} & \hd{리허설 2} & \hd{리허설 3} & \hd{본 컷오버 마스터} & \hd{본 컷오버 이력} & \hd{서명} \\ \hline
\ifwbblank
\rd{8mm}① 건수 & \g{[n테이블 행 수]} & \g{[원본 $-$ 적재]} & \g{0} & \g{[결과·실패 항목]} & & & & & \\ \hline
\rd{8mm}② 금액 합계 & \g{[금액 원장 n년]} & \g{[원본 합 $-$ 적재 합, 원 단위]} & \g{0} & & & & & & \\ \hline
\rd{8mm}③ 키 유일성 & \g{[PK n]} & \g{[중복 PK 건수]} & \g{0} & & & & & & \\ \hline
\rd{8mm}④ 참조 정합성 & \g{[FK n]} & \g{[고아 레코드 건수]} & \g{0} & & & & & & \\ \hline
\rd{8mm}⑤ 코드 매핑 커버리지 & \g{[매핑표 n 코드]} & \g{[미매핑 ÷ n]} & \g{0 (100\%)} & & & & & & \\ \hline
\rd{8mm}⑥ 표본 필드 대조 & \g{[무작위 n행 × 전 필드]} & \g{[불일치 행 수]} & \g{0 (재실행 n회)} & & & & & & \\ \hline
\else
① 건수 & 1,842테이블 행 수 & 원본(필터 후) 행 수 $-$ 적재 행 수 & 0 & 0 (7.4TB 전체) & 0 & 0 & 0 (2,184,310) & 0 (3.21억) & 용역사 PM·P1 PMO 품질 \\ \hline
② 금액 합계 & 매출·발주·재고 금액 원장 5년 & 원본 합계 $-$ 적재 합계 (원 단위) & 0 & $-$41,200원(반올림) → 규칙 정정 & 0 & 0 & 0 & 0 (5년 매출 합 일치) & 재경팀장 \\ \hline
③ 키 유일성 & PK 1,842 & 중복 PK 건수 & 0 & \textbf{1,204}(상품 코드 브랜드 간 중복 — D-0958·R-10) & 0 (정본 정제 후) & 0 & 0 & 0 & 박세연 \\ \hline
④ 참조 정합성 & FK 3,916 & 고아 레코드 건수 & 0 & \textbf{412}(폐점 매장·단종 상품 참조) & 0 (매핑 보정) & 0 & 0 & 0 & 박세연 \\ \hline
⑤ 코드 매핑 커버리지 & 코드 매핑표 4,318 코드 & 미매핑 코드 수 ÷ 4,318 & 0 (100\%) & \textbf{96}(97.8\%) & 0 & 0 & 0 & 0 & 수행 통합팀·P1 PMO \\ \hline
⑥ 표본 필드 대조 & 무작위 표본 2,000행 × 전 필드(마스터) / 5,000행(이력) & 필드값 불일치 행 수 & 0 (재실행 2회 허용) & 18 & 3 (날짜 시간대·통화 표기) & 0 & \textbf{1}(가격 마스터 VAT 포함가 소수 자릿수) → 재실행 → 0 & 0 & P1 PMO 품질·용역사 PM \\ \hline
\fi
\end{tabularx}}
\B{\small\g{리허설 1회차 실패의 원인과 해소 경로(정본 판정·필터·규칙) / 본 컷오버 불일치의 원인: }\hrulefill\par}{\small 리허설 1회차의 창 초과(41h 20m)는 7.4TB 전체를 대상으로 했기 때문이며 CR-11(5년 필터)이 G2 정본 판정(11-27)과 함께 확정된 뒤 2회차부터 5.1TB로 실행했다. 1회차 미통과 3종은 모두 \textbf{데이터 정본 정제}(G2)로 해소되었고, 본 컷오버의 표본 1건은 변환 규칙(반올림)의 문제였다 — 리허설 3회차가 통과했는데 본 컷오버에서 나온 이유는 2월 가격 개정분(신규 가격 12건)이 3회차 이후 추가되었기 때문이다.\par}

% ------------------------------------------------------------------ 7 2차 컷오버 축약 런북
\secbar[110mm]{7. 2차 컷오버 축약 런북 — 두 컷오버 사이의 임시 상태 \B{}{(물류: 이천센터 재고 원장 · P2 설비 인터페이스 4본, v1.0 03-19 확정)}}
\guide{2차 창, 착수 조건(\textbf{1차 하이퍼케어 지표}·선행 프로젝트 MC·연동 시험), 축약 시각표(10행 안팎), 게이트 3개의 축약판, 롤백 트리거 2개, 그리고 \textbf{두 컷오버 사이의 임시 상태} — 어느 원장이 정본인가, 그 상태에서 하지 말아야 할 것, 임시 연동의 폐기 시점. 2차 컷오버를 “1차와 동일”로 적지 않는다. 2차 go 조건이 1차 하이퍼케어와 연결되어 있지 않으면 R-14 트리거가 문서 밖에 남는다.}
{\footnotesize
\begin{tabularx}{\linewidth}{|L{30mm}|Y|}\hline
\hd{항목} & \hd{내용} \\ \hline
\B{\rd{8mm}}{}창 & \Bg{[YYYY-MM-DD hh:mm ~ hh:mm, n시간 — 시간대 선택 근거, SLA 산정 처리]}{2027-03-27(토) 22:00 ~ 03-28(일) 06:00, 8시간(센터 야간 출고 없는 시간대 [추가 설정]). SLA 산정 제외(계획 정지)} \\ \hline
\B{\rd{16mm}}{}착수 조건 (게이트 1) & \Bg{[① 1차 하이퍼케어 n주 지표 — 실제 ② 선행 MC·연동 시험 ③ 업무 성수 창 회피 확인(실명) ④ 실사 등 사전 작업]}{(03-24 운영위 + 03-27 21:30) ① \textbf{1차 하이퍼케어 2주(03-01~12) 심각도 2 이상 5건 이하 — 실제 4건}(R-14 트리거 미발동) ② P2 MC 03-05 확인, A15 P2 인터페이스 시험 03-08~12 통과 ③ 3월 B2B 납품 집중 창(설 이후) 회피 확인 — 오정근 03-19 ④ 센터 재고 실사(순환) 03-26 완료 — 한동석} \\ \hline
\B{\rd{20mm}}{}임시 상태 (두 컷오버 사이) & \Bg{[기간 / 어느 원장이 정본인가(매장·센터) / 대사 방법 / \textbf{하지 말 것} / 리포트의 분모 처리 / 임시 연동 폐기 시점]}{02-28 ~ 03-28(4주). 매장 재고 = 신 원장(정본), 이천센터 재고 = 구 SAP 원장(정본), 둘 사이 대사는 폴백 야간 배치 2회(21:00·03:00). \textbf{하지 말 것}: 센터 재고를 신 원장에서 조회·수정하지 않는다(읽기 전용 미러), 재고 정확도 리포트(StRS-034)는 3월분을 “매장 원장 기준”으로만 산출하고 센터 분모는 4월부터. CR-08 임시 연동 2본 폐기는 2차 후 30일 = \textbf{04-27}(CCB 01-21 갱신)} \\ \hline
\B{\rd{22mm}}{}축약 시각표 & \Bg{[hh:mm 게이트 1 → 정지·마감 → 추출·적재 → 검증 n종 → 게이트 2(판정자) → PoNR → 연동 활성·시험 → 첫 배치 → 게이트 3·오픈 판정(판정자) → 오픈]}{21:30 게이트 1 → 22:00 센터 출고 정지·구 원장 마감 → 22:30 센터 재고 추출·적재(84,000 SKU-로케이션) → 00:30 검증 3종(건수·금액·표본 500) → \textbf{01:30 게이트 2}(한동석·박세연) → 01:40 센터 원장 신 플랫폼 활성화(PoNR) → 02:00 P2 설비 IF 4본 활성·수신 시험 → 03:00 폴백 배치 1회차(통합 원장 기준) → \textbf{04:30 게이트 3·오픈 판정}(정미란 위임 박세연) → 06:00 센터 가동} \\ \hline
\B{\rd{12mm}}{}트리거 (2개) & \Bg{[RB-A 지표·임계값·판정자·조치 / RB-B — 부분 모드 오픈 가능 여부]}{RB-A 센터 재고 대사 불일치 0.5\% 초과(게이트 2, 한동석 → 롤백, 손실 0) / RB-B 설비 IF 4본 중 1본 이상 수신 실패(02:00~04:00, 박세연 → 설비 IF만 수동 모드로 오픈, 센터 원장은 유지)} \\ \hline
\B{\rd{12mm}}{}실행 기록 & \Bg{[게이트별 실제 시각·판정, 트리거 관측, 오픈 시각, 2차 후 첫 평일에 일어난 일]}{게이트 2 01:25 go(대사 불일치 0.03\% = 25건, 실사 미반영 → 수동 보정), PoNR 01:40, 설비 IF 4/4, 오픈 06:00. 트리거 발동 0. \textbf{03-29(월) 03:00 폴백 배치 1회차 실패(P2-07, 84분 지연 — KE-02)} — 2차 후 첫 평일 배치에서 발생, 22:30 파라미터로 조정} \\ \hline
\end{tabularx}}

% ------------------------------------------------------------------ 8 SAC
\secbar[120mm]{8. 서비스 수용 기준 (SAC) \B{v[ ] → 판정}{v1.0 → 판정 — 초안 G1 2026-04-24 합의, 월 검토 9회, 사전 점검 2027-01-28, 판정 2027-05-21, 판정권자 박준영 (RACI 22행)}}
\guide{10~12항목. 운영 조직의 관심 순서로 — 문서·모니터링·백업복구·보안·교육·KEDB·SLA 서명·절차·소스 인수·인터페이스 운영·계약 이관·잔여 결함. 조건은 관측 가능하게, \textbf{판정자는 인수자}(운영 조직 — 기술 항목은 기술 인수자, 보안은 CPO, 계약은 구매). 증빙은 문서·기록(“완료 보고”가 아니다). G3 사전 점검 상태와 G4 판정 상태를 나란히 — 사전 점검이 없으면 G4에서 처음 본다. 미해소 시 처리 = 해소 기한·책임자 서명 → 종료 보고서의 미해결 사항 이관 목록.}
{\footnotesize\setlength{\tabcolsep}{2pt}
\begin{tabularx}{\linewidth}{|C{6mm}|L{24mm}|Y|L{22mm}|L{30mm}|L{28mm}|L{46mm}|L{22mm}|}\hline
\hd{\#} & \hd{항목} & \hd{조건 (관측 가능)} & \hd{판정자} & \hd{증빙} & \hd{G3 사전 점검 \B{[일자]}{01-28}} & \hd{G4 판정 \B{[일자]}{05-21}} & \hd{미해소 시} \\ \hline
\ifwbblank
\rd{9mm}1 & \g{[운영 문서 n종]} & \g{[n종 전부 v1.0 이상 인수, 운영팀 검토 서명]} & \g{[인수자]} & \g{[문서 목록·서명표]} & \g{[미해소 (n/n)]} & \g{[충족 / 조건부 / 미해소]} & \g{[기한·책임자·서명]} \\ \hline
\rd{9mm}2 & \g{[모니터링·알림]} & \g{[임계값 n개, 알림 경로 시험 n회]} & & & & & \\ \hline
\rd{9mm}3 & \g{[백업·복구]} & \g{[RPO·RTO, DR 전환 시험 실측]} & & & & & \\ \hline
\rd{9mm}4 & \g{[보안]} & \g{[진단 High 0, 권한 매트릭스, 파기 증적]} & \g{[CPO]} & & & & \\ \hline
\rd{9mm}5 & \g{[교육]} & \g{[집단별 이수율]} & & & & & \\ \hline
\rd{9mm}6 & \g{[Known Error DB]} & \g{[ELS 중 식별된 KE 전부 등재, 우회 방법]} & & & \g{[해당 없음(ELS 전)]} & & \\ \hline
\rd{9mm}7 & \g{[SLA/OLA/UC 서명]} & \g{[SLA n·OLA n·UC n 전부 서명 — 한 층 미서명이면 발효 불가]} & & & & & \\ \hline
\rd{9mm}8 & \g{[인시던트·문제·변경 절차]} & \g{[우선순위 정의, CAB 운영 n회 이상]} & & & & & \\ \hline
\rd{9mm}9 & \g{[소스·형상 인수]} & \g{[클린 빌드 재현 100\%, SBOM, 절차서]} & \g{[기술 인수자]} & & & & \\ \hline
\rd{9mm}10 & \g{[인터페이스 운영]} & \g{[n본 문서, 배치 절차, 임시 연동 폐기 확인]} & & & & & \\ \hline
\rd{9mm}11 & \g{[라이선스·계약 이관]} & \g{[계약별 이관 확인, 절차 승인]} & \g{[구매]} & & & & \\ \hline
\rd{9mm}12 & \g{[잔여 결함·ELS 종료]} & \g{[심각도 1·2 미해결 0, 잔여 목록 서명, ELS 종료 기준 충족]} & & & \g{[해당 없음]} & & \\ \hline
\else
1 & 운영 문서 12종 & 12종 전부 v1.0 이상 인수, 운영팀 검토 서명(⑰ 항목 2) & 박준영 & 문서 목록·서명표 & 미해소 (8/12) & \textbf{충족} (12/12, 10번 아카이브 절차는 문서 있음·승인 대기 → 항목 11) & — \\ \hline
2 & 모니터링·알림 & APM 대시보드 임계값 정의 24개, 알림 경로(서비스데스크·당직) 시험 3회 & 박준영 & 임계값 정의서, 알림 시험 기록 & 미해소 (임계값 12개) & 충족 (24개, ELS 중 조정 이력 포함) & — \\ \hline
3 & 백업·복구 & RPO 15분·RTO 4시간 [추가 설정, P5 초안], DR 전환 시험 1회 실측 & 박준영·박세연 & DR 전환 시험 보고(04-10, RTO 3h 20m) & 미해소 (시험 미실시) & 충족 & — \\ \hline
4 & 보안 & 보안진단 2회차 High 0(01-20), 계정·권한 매트릭스 승인, 개인정보 파기 배치 증적 3회 & 최영주 & 진단 보고, 권한 매트릭스, 파기 로그 & 충족 & 충족 (03-05 사건 후 이벤트 중복 검출 추가) & — \\ \hline
5 & 교육 & 서비스데스크 11명·운영팀 8명 교육 이수 100\%, 직영 매장 관리자 227점 100\%, 가맹 점주 84\%(P3) & 박준영·오정근 & 이수 대장 & 미해소 (데스크 6/11) & 충족 (가맹 84\%는 P3 목표 80\% 이상) & — \\ \hline
6 & Known Error DB & ELS 중 식별된 Known Error 전부 등재(우회 방법·근본 원인 상태), 인시던트–문제 연결 & 박준영 & KEDB 6건(⑰ 항목 6) & 해당 없음(ELS 전) & 충족 & — \\ \hline
7 & SLA/OLA/UC 서명 & SLA 1(현업↔IT), OLA 3, UC 5 전부 서명 — 한 층이라도 미서명이면 SLA 발효 불가 & 박준영·박세연 & 서명본 9건(⑰ 항목 4) & 미해소 (UC 3PL 없음, 수행사 유지보수 UC 초안) & 충족 (UC 3PL은 “배치 수신 확인”으로 한정 서명, SLA 출고 항목 최선 노력으로 강등) & — \\ \hline
8 & 인시던트·문제·변경 절차 & 인시던트 우선순위 P1~P4 정의, 문제 관리 주기, CAB 수요일 운영 4회 이상 실시 & 박준영 & 절차서, CAB 회의록 & 미해소 (CAB 미실시) & 충족 (CAB 03-17·24·31·04-07… 8회) & — \\ \hline
9 & 소스·형상 인수 & 발주사 저장소 클린 빌드 재현 100\%(하도급 220 FP 포함), SBOM·라이선스 목록, 빌드·배포 절차서 & 박세연·형상 담당 & 재현 시험 기록(05-12) & 충족 (일 1회 동기화, 월 1회 재현) & 충족 & — \\ \hline
10 & 인터페이스 운영 & 47본 설계·운영 문서(AF-1-02 9본 보완 완료), 폴백 배치 운영 절차(22:30 파라미터), 임시 연동 2본 폐기 확인(04-27) & 박세연·박준영 & 문서 47/47, 폐기 확인서 & 미해소 (폐기 전) & 충족 & — \\ \hline
11 & 라이선스·계약 이관 & 상용SW 유지보수(P5, SW-04)·클라우드(CL-01)·앱 벤더(앱-01)·API GW 정식 라이선스(03-05) 이관, \textbf{아카이브 접근 절차 재경 승인}(CR-11 조건 ④) & 강현도·박준영 & 이관 확인서, 재경 승인서 & 미해소 & \textbf{조건부 충족} — 아카이브 접근 절차 재경 승인 대기(기한 06-15, 책임자 재경팀장·박준영 서명) → ⑱ 미해결 사항 이관 1번 & 기한·책임자 서명 \\ \hline
12 & 잔여 결함·ELS 종료 & 심각도 1·2 미해결 0, 심각도 3 이하 잔여 목록 서명, \textbf{ELS 종료 기준 충족(04-23)} & 박준영 & ELS 종료 보고, 잔여 결함 목록(심각도 3: 11, 4: 27) & 해당 없음 & 충족 & — \\ \hline
\fi
\end{tabularx}}
\B{\small\g{판정: n항목 중 충족 n · 조건부 n → 헌장 완료 기준 ① 판정 / 서명(인수자·일자) / 승인(스폰서·일자): }\hrulefill\par}{\small 판정: \textbf{12항목 중 충족 11·조건부 1 → 헌장 §14 완료 기준 ① 충족}(미해소 0 또는 해소 기한·책임자 서명). 서명 박준영 05-21, 승인 정미란 05-28(G4).\par}
\landscapeoff

% ------------------------------------------------------------------ 9 컷오버 직전 사건 기록 (세로)
\secbar[110mm]{9. 컷오버 직전 사건 기록 \B{}{— D$-$2 API 게이트웨이 라이선스}}
\guide{창 앞에 발견된 사건의 발견 경위·시각, 사실, 영향, 선택지(각각의 결과), 결정·결정자·시각, 비용과 자금원, 사후 조치, \textbf{원인}(“담당자 실수”가 아니라 절차·양식이 만든 대기 — 절차가 만든 대기를 사람 탓으로 적으면 교훈이 남지 않는다), 교훈 등록부 연결. 사건이 게이트 1 조건에 어떻게 반영되었는지를 적는다.}
{\footnotesize
\begin{tabularx}{\linewidth}{|L{22mm}|Y|}\hline
\hd{항목} & \hd{내용} \\ \hline
\B{\rd{10mm}}{}발견 & \Bg{[YYYY-MM-DD hh:mm, 어느 점검에서 누가 무엇을 조회해 무엇을 알았는가]}{2027-02-24(수) 14:30, 컨트롤룸 사전 점검(행 3 준비) 중 P1 PMO 기준선 담당이 운영 환경 라이선스 키 발급 상태를 벤더에 조회 → “발주서 미접수” 회신} \\ \hline
\B{\rd{16mm}}{}사실 & \Bg{[계약·금액·발주 단계·어느 결재에서 언제부터 대기 — 지침 조항·시행일 인용]}{API GW 라이선스 5.6억(SW-01) 중 운영 환경 2차분(개발·시험 환경 1차분은 2026-03 인도) 발주서가 02-13 소관 본부장(박세연) 결재 후 \textbf{CFO 직접 결재 대기}(재무본부 내부통제 지침 2026-02-16 시행: 단일 건 5,000만 원 이상) — 정미란 G3 준비·해외 출장(02-16~23)으로 결재함 정체. 구매팀은 “결재 완료 후 발주”가 절차이므로 벤더에 발주서 미발송} \\ \hline
\B{\rd{8mm}}{}영향 & \Bg{[컷오버에 미치는 영향 — 창 이탈 여부]}{정식 키 없이 02-26 컷오버 시 운영 환경 API GW 기동 불가(개발 키는 계약상 운영 사용 금지) → 창 이탈} \\ \hline
\B{\rd{14mm}}{}선택지 & \Bg{[① 대기(소요·창 영향) ② 임시 수단(비용·조건) ③ 기각 사유]}{① 결재 완료 대기(정미란 02-25 복귀, 결재 02-25 → 발주 → 벤더 키 발급 표준 5영업일 → 03-04) — 창 이탈 ② \textbf{벤더 임시 라이선스 30일}(수수료 0.02억, 정식 발주 시 상계 불가) ③ 개발 키 운영 사용 — 계약 위반, 기각} \\ \hline
\B{\rd{12mm}}{}결정 & \Bg{[채택안 — 결정자·시각, 적용·시험 시각, 자금원(풀·보고 일정), 런북 반영]}{\textbf{② 채택} — P1 PM·강현도 02-24 17:00, 벤더 담당 회신 02-24 19:00, 임시 키 02-25 10:00 적용·시험. 자금원: 컨틴전시 미배정 풀 0.02(CCB 03-18 사후 보고). 게이트 1 조건 ③에 “임시 라이선스 유효” 추가(런북 v1.1)} \\ \hline
\B{\rd{8mm}}{}사후 & \Bg{[정식 결재·발주·적용 일자, 만료 전 확인자]}{정미란 결재 02-25 → 발주 02-26 → 정식 키 \textbf{03-05 적용}(임시 키 만료 03-26 전). 재경팀장 확인} \\ \hline
\B{\rd{14mm}}{}원인 & \Bg{[절차의 부작용 — 어느 지침에 무엇이 없었는가, 어느 양식이 무엇을 추적하지 않았는가]}{담당자 실수가 아니라 \textbf{절차의 부작용} — 프로젝트 컷오버 창 앞 2주 동안 5,000만 원 이상 발주 목록을 사전에 결재해 두는 규칙이 지침에 없음. G3 자료의 “조달 상태” 절(⑮ 항목 8)에 라이선스 2차분이 “계약 내”로 표기되어 발주 상태가 추적되지 않음} \\ \hline
\B{\rd{8mm}}{}교훈 & \Bg{[L-nn — 개정 요청 대상 표준·양식]}{\textbf{L-09}(⑲): 내부통제 지침 개정 요청 — “컷오버 창 앞 사전 결재 목록” 조항; G3 조달 상태 표에 “발주·결재 단계” 열 추가} \\ \hline
\end{tabularx}}

% ------------------------------------------------------------------ 10 실행 기록 · 인수 서명
\secbar[60mm]{10. 실행 기록 · 인수 서명}
\guide{게이트 판정 기록(시각·판정·서명자), 트리거 관측(발동 0이어도 관측값/임계값), 계획 대비 편차(분 단위 — “약간 지연”은 기록이 아니다), D+1 ELS 착수 인계 서명(인수자 실명·시각), 2차 컷오버 인수. 성공한 컷오버의 기록을 남기지 않는 것이 가장 흔한 오류다.}
{\footnotesize
\begin{tabularx}{\linewidth}{|L{28mm}|Y|}\hline
\hd{항목} & \hd{내용} \\ \hline
\B{\rd{12mm}}{}게이트 판정 기록 & \Bg{[1: 일시 go/no-go(서명) / 2: (재실행 n회) / 3: / 오픈:]}{1: 02-26 19:35 go(P1 PM·이재현·박준영) / 2: 02-27 01:38 go(박세연·용역사 PM·P1 PM — 재실행 1회) / 3: 02-27 12:30 go(정미란·박세연·P1 PM) / 오픈: 02-28 04:30 go(정미란·박준영·이재현·P1 PM)} \\ \hline
\B{\rd{10mm}}{}트리거 관측 & \Bg{[RB-1 관측값/임계값 · RB-2 · RB-3 · RB-4 → 발동 n]}{RB-1 재실행 1/2 · RB-2 08:35/11:00 · RB-3 9/31(직영 2/10) · RB-4 100\%/99.0\%, PG 0/0 → \textbf{발동 0}} \\ \hline
\B{\rd{10mm}}{}편차 & \Bg{[게이트별 ±분, 적재, 배포, 오픈 / 여유 소진 n / n]}{게이트 2 +38분(표본 대조), 게이트 3 +30분(스폰서 열람), 적재 $-$25분, 배포 +10분, 오픈 0. 여유 소진 1h 10m / 2h 20m} \\ \hline
\B{\rd{12mm}}{}D+1 ELS 착수 인계 & \Bg{[인수자 실명 · 서명 일시 · ELS 착수 주차 · 컨트롤룸 해산·워룸 규모]}{\textbf{박준영 2027-03-01 07:20 서명} — ELS 1주차 착수(⑰ 항목 5), 컨트롤룸 해산·워룸 축소(수행사 20 → 15)} \\ \hline
\B{\rd{10mm}}{}2차 컷오버 & \Bg{[오픈 일시, 인수자 서명, 트리거 발동]}{03-28 06:00 오픈, 한동석 인수 서명 03-28 07:00, 트리거 발동 0} \\ \hline
\end{tabularx}}
\vspace{3mm}
\noindent{\footnotesize
\begin{tabularx}{\linewidth}{|Y|Y|Y|Y|}\hline
\hd{통제관 (서명)} & \hd{게이트 3 · 오픈 판정자 (서명)} & \hd{운영 인수자 (서명)} & \hd{수행사 PM (서명)} \\ \hline
\Bg{[P1 PM]}{P1 PM\rd{10mm}} & \Bg{[스폰서]}{정미란 02-28 04:30} & \Bg{[서비스운영팀장]}{박준영 03-01 07:20} & \Bg{[수행사]}{이재현 02-28 04:30} \\ \hline
\end{tabularx}}
\end{document}
